% ============================================================
% RESUMEN BILINGUE - 200 palabras por idioma (Roselyn Sanchez)
% ============================================================
\section*{Resumen}
\addcontentsline{toc}{section}{Resumen}

Este informe consolida el cierre del Proyecto Integrador de Ingenier\'ia de Requisitos: la especificaci\'on del Sistema de Gesti\'on Agr\'icola con Inteligencia Artificial (SGA) para una finca de verde y cacao de 16 hect\'areas en AgroMoreira, cuyo registro manual y comunicaci\'on verbal generaban p\'erdida de informaci\'on, proyecciones de producci\'on err\'oneas y reacci\'on tard\'ia ante enfermedades como moniliasis y sigatoka. Sobre la l\'inea base auditada se consolid\'o la ERS v2.0 con 57 requisitos (31 funcionales y 26 no funcionales), dos componentes de inteligencia artificial especificados con umbrales verificables ---un recomendador de manejo agr\'icola (IA-01) y un detector de plagas por imagen (IA-02)---, clasificaci\'on de riesgo conforme al Reglamento (UE) 2024/1689 y cumplimiento de la Ley Org\'anica de Protecci\'on de Datos Personales del Ecuador. La auditor\'ia de calidad midi\'o seis m\'etricas derivadas de ISO/IEC 25010:2023 con aritm\'etica visible; las m\'etricas corregidas pasaron de 3,50 a 3,00 (modificabilidad) y de 0,119 a 0,024 defectos por requisito (correcci\'on). La trazabilidad end-to-end enlaz\'o evidencia de campo, requisitos, casos de uso, clases, criterios BDD y casos de prueba conceptuales. El documento incluye retrospectiva del equipo, conclusiones por unidad, declaraci\'on de uso de IA y sustentos para la defensa individual ante el tribunal.

\vspace{0.4cm}
\noindent\textbf{Palabras clave:} ingenier\'ia de requisitos; ERS/SRS; trazabilidad; m\'etricas de calidad ISO/IEC 25010; requisitos de inteligencia artificial; agroindustria.

\clearpage
\section*{Abstract}
\addcontentsline{toc}{section}{Abstract}

This report closes the Requirements Engineering integrative project: the specification of an Artificial Intelligence Agricultural Management System (SGA) for a 16-hectare green plantain and cacao farm in AgroMoreira, where manual record keeping and verbal communication caused information loss, inaccurate production forecasts and late reaction to diseases such as black pod (moniliasis) and sigatoka. Building on the audited baseline, the consolidated SRS v2.0 delivers 57 requirements (31 functional and 26 non-functional) and two artificial intelligence components specified with verifiable thresholds ---an agricultural management recommendation engine (IA-01) and an image-based pest detector (IA-02)--- including risk classification under Regulation (EU) 2024/1689 and compliance with Ecuador's Organic Personal Data Protection Law. The quality audit measured six metrics derived from ISO/IEC 25010:2023 with visible arithmetic; corrected metrics improved from 3.50 to 3.00 (modifiability) and from 0.119 to 0.024 defects per requirement (correction). End-to-end traceability linked field evidence, requirements, use cases, classes, BDD acceptance criteria and conceptual test cases. The report includes a team retrospective, unit-by-unit conclusions, an AI usage declaration and supporting material for the individual technical defense before the academic panel.

\vspace{0.4cm}
\noindent\textbf{Keywords:} requirements engineering; SRS; traceability; ISO/IEC 25010 quality metrics; artificial intelligence requirements; agribusiness.
